\subsection{NN-G035: a complete first-cell mass gate for competitive closed tables}
\paragraph{Statement.}
Every symmetric closed-table effective associative lift with
at most $2744$ cancelling pairs lies in the complete necessary
mass relaxation below. It has $w\in\{8,10,12\}$ and $s\ge588$.
Its first-cell mass triple $(g_1,d_1,u_1)$ belongs to the
following set of $21$ triples, grouped by $g_1$:
\[
\begin{array}{c|l}
 g_1&\text{possible }(d_1,u_1)\\\hline
 64&(4,4)\\
 72&(4,6),(6,4)\\
 76&(4,6)\\
 80&(4,8),(8,4),(8,6),(12,4)\\
 84&(6,6),(10,6),(14,4)\\
 88&(8,8)\\
 96&(4,10),(4,12),(6,6),(12,4),(14,6)\\
 100&(4,10),(6,10),(10,6)\\
 128&(12,8)
\end{array}                                                         \tag{G035.1}
\]
The full relaxation has $43$ base rows and $467$ possible
auxiliary increment pairs, all explicitly listed in the
certificate. Its minimum relaxed score is $1700$.
Consequently the minimum over all symmetric closed tables
is between $1700$ and $2744$.
Surviving mass tuples are not asserted realizable; no
Laurent degree or support-location bound is inferred.
This theorem does not apply to the broader unrestricted
$1000$-position problem of NN-G029.

\paragraph{Complete finite bounds.}
Use the actual-coefficient mass notation of (G031.1).
The known lower bounds are
\[
 s\ge576,\quad g_i\ge48,\quad d_i\ge4,\quad
 u_i\ge4,\quad v_i\ge16,\quad a_1\ge8,\quad a_2\ge16.
\]
All are even; NN-G030 also makes $g_1,g_2,s,a_2$
multiples of four. The score cap is $T\le9912$.
NN-G031 gives $w\ge8$, and the mass formula implies
\[
              20(g_1+g_2)\le9912-3520,
              \qquad g_1+g_2\le316.                 \tag{G035.2}
\]
Thus each $g_i\le268$.
If $w\ge20$, the relation $g_1g_2=ws\ge576w$
and the arithmetic-geometric mean inequality give
$g_1+g_2\ge48\sqrt{20}>214$.
Since the sum is divisible by four, it is at least
$216$. The mass formula would then imply
\[
 T\ge576+20\cdot216+200\cdot20+800+208+336
       =10240>9912.
\]
Hence $w$ is one of $8,10,12,14,16,18$.

Fix one such $w$ and multiples of four $g_i\ge48$
with sum at most $316$. Then $s=g_1g_2/w$ must be
an integer multiple of four at least $576$.
For each index $i$, with $j\ne i$, enumerate the
even integers $4\le d_i\le\lfloor g_i/4\rfloor$,
retaining exactly those for which
\[
             u_i=\frac{w(g_i-4d_i)}{g_j}
\]
is an even integer at least four. These finite bounds
contain every possible actual cell mass.

The auxiliary masses satisfy
$v_i=w(4g_i-a_i)/g_j$. Set $t_1=2,t_2=4$, and
\[
 M_i=\operatorname{lcm}\left(t_i,
                   \frac{2g_j}{\gcd(w,2g_j)}\right).
\]
Their necessary congruence is exactly
$a_i\equiv4g_i\pmod{M_i}$, together with
$a_1\ge8,a_2\ge16$.
Let $a_i^0$ be the least integer satisfying this lower
bound and congruence, and put
$v_i^0=w(4g_i-a_i^0)/g_j$. If $v_i^0<16$, discard
the tuple, since increasing $a_i$ only decreases $v_i$.
Otherwise every actual auxiliary mass is
\[
 a_i=a_i^0+M_i k_i,\qquad
 v_i=v_i^0-\frac{wM_i}{g_j}k_i,\qquad k_i\ge0.       \tag{G035.3}
\]
The mass increment per unit $k_i$ is
\[
 \Delta_i=14M_i-20\frac{wM_i}{g_j}>0,
\]
because $w\le18$ and $g_j\ge48$ make
$14-20w/g_j\ge13/2>0$.
Thus the least auxiliary masses give a rigorous lower
bound $T_0$ for the tuple, and all remaining choices
satisfy the explicit finite inequalities
\[
 \Delta_1k_1+\Delta_2k_2\le9912-T_0,\qquad
 \frac{wM_i}{g_j}k_i\le v_i^0-16.                   \tag{G035.4}
\]

\paragraph{The extra $U_i=D$ tests for masses eight and ten.}
For $w=8,10$, (G034.5) shows that $W$ is primitive
and coprime to $D$.
If $u_1=4$, then $U_1=D$. The proof of (G031.4)
gives $G_2=DWC$ for an integral $C\equiv D-1\pmod2$;
therefore $g_2/(4w)$ is an odd integer.
If $u_2=4$, the analogous quotient is even, so
$8w\mid g_1$. These tests use only the stated mass
of $W$ and integral associativity, not a constant-ratio
or a degree assumption.

\paragraph{Why the minimal bottom polynomial is impossible.}
Suppose $s=576$. Coefficientwise domination and equality
of masses force
\[
                  S=4D^2(D-1)^2.                    \tag{G035.5}
\]
Write $C=D-1$ temporarily. From $WS=G_1G_2$ and
$G_1\equiv0,G_2\equiv D^2C\pmod2$, reduction after
dividing by two shows that $G_1$ is divisible by four.
Writing $G_1=4J$ gives
\[
             JG_2=WD^2C^2,\qquad J\equiv DC\pmod2.
\]
Since $W\equiv D\pmod2$, $W(1)$ is odd, so $C$
does not divide $W$. Neither $J$ nor $G_2$ can be
divisible by $C^2$, by their displayed reductions
modulo two. Their product contains exactly $C^2$;
therefore each contains exactly one factor $C$.
Write
\[
          J=CA,\qquad G_2=CB,\qquad
          AB=WD^2,\qquad A\equiv D, B\equiv D^2\pmod2.
\]
The actual $D$-adic orders of $A,B,W$ are at most
$1,2,1$, respectively, again by these nonzero reductions.
If $D\nmid W$, their first two orders sum to two;
if $D\mid W$, they sum to three. In either case
at least one is maximal: either $A=DA_1$ with
$A_1\equiv1\pmod2$, or $B=D^2B_1$ with
$B_1\equiv1\pmod2$.
Thus at least one of $g_1,g_2$ is $48$ times a
positive odd integer.

It cannot equal $48$: using $g_1g_2=576w$ in
the corresponding $u_i$ identity gives
$u_i=4-d_i/3<4$.
Hence some $g_i\ge144$.
For that index, eliminate $v_i$ using $s=576$:
\[
 20v_i+14a_i=\frac{5g_i^2}{36}
                   +\left(14-\frac{5g_i}{144}\right)a_i.
\]
The second coefficient is positive because $g_i\le268$.
Since $a_i\ge8$, the right side is at least
$5g_i^2/36+112-5g_i/18$, an increasing function for
$g_i\ge144$, and hence at least $2952$.
For the other index $20v_j+14a_j\ge432$.
Also $g_1+g_2\ge192$ and $w=g_1g_2/576\ge12$.
Together these estimates imply
\[
 T\ge576+3840+2400+160+208+2952+432
       =10568>9912.
\]
This contradiction excludes every minimal-bottom case;
no finite computation was used for this polynomial step.

\paragraph{Exact enumeration and completeness of the certificate.}
The verifier enumerates precisely the finite ranges in
(G035.2)--(G035.4), using integer gcd and lcm arithmetic.
It first tests $s$, the lower mass bound, and the two
$U_i=D$ divisibility tests, then enumerates every $d_i$
and uses the least admissible $a_i^0$.
There are $1445$ coarse $(w,g_1,g_2)$ tuples before
the cell tests and $55$ surviving base rows with
$T_0\le9912$. Exactly $12$ have $s=576$ and are
excluded by the preceding proof. The remaining $43$
rows have exactly $467$ increment pairs (G035.4).
The certificate lists every row and every pair.
It explicitly includes the actual NN-G033 mass tuple.

Projecting these rows gives (G035.1), $w\in\{8,10,12\}$,
and least $s=588$. Their least $T_0$ is $7824$, giving
the relaxed lower bound $(7824-4424)/2=1700$.
Every actual realization at or below the incumbent
maps to one of these rows and increments by the
proved equations and bounds. Larger scores are already
above $1700$. This proves the stated closed-table
lower bound and all finite necessary restrictions.
\hfill$\square$

\paragraph{Geometric translation and scope.}
For the actual graded sign-representation kernels of
NN-G010, this theorem bounds the possible ordinary
dimensions of the first-cell coefficient object, its
bottom action object, and its correction product by
the $21$ listed triples. Their grading locations are
still unrestricted. The bottom coefficient object
cannot have exactly the absolute trace character
$4D^2(D-1)^2$ in any competitive realization; integral
character factorization and parity force a dimension
contradiction. The finite list records necessary
dimension data, not actual objects for every row.

The dictionary is SGA 3, I, 4.7.2--4.7.3 and archive
cote 2, scan 17 via NN-G010. No equality of formal
$K_0$ classes is treated as an effective realization.
The result shrinks the closed-table search to explicit
mass data, while NN-G029--NN-G032 separately address
the wider unrestricted presentation. Neither statement
supplies a canonical GCT basis, rank-four positivity,
or PH0/PH1.

